\chapter{Kompaktlık}

Kompaktlık, sonsuz yapılarda ``sınırlılık'' fikrinin topolojik soyutlamasıdır.
Açık örtü tanımı temel alınarak incelenmekte; varyantlar ve lokal kompaktlık da ele alınmaktadır.

%-------------------------------------------------------------------
\section{Konu}

\subsection{Kompaktlık Tanımı}

\begin{definition}[Kompakt Uzay]
$(X, \tau)$ topolojik uzayı \textbf{kompakt} ise:
Her açık örtü $\{U_\alpha\}$ için sonlu alt-örtü vardır:
\[
X \subseteq \bigcup_{\alpha} U_\alpha
\implies
\exists \text{ sonlu } I:\; X \subseteq \bigcup_{\alpha \in I} U_\alpha.
\]
\end{definition}

\subsection{Kompaktlık Varyantları}

\begin{table}[h]
\centering
\begin{tabular}{ll}
\toprule
\textbf{Varyant} & \textbf{Tanım} \\
\midrule
Kompakt & Her açık örtünün sonlu alt-örtüsü \\
Sayılabilir kompakt & Her sayılabilir örtünün sonlu alt-örtüsü \\
Ardışımlı kompakt & Her dizi yakınsak alt-dizi içerir \\
Sözde kompakt & Her sürekli $f: X \to \mathbb{R}$ sınırlıdır \\
Lindelöf & Her açık örtünün sayılabilir alt-örtüsü \\
$\sigma$-kompakt & Sayılabilir kompakt kümelerin birleşimi \\
Lokal kompakt & Her noktanın kompakt komşuluğu var \\
\bottomrule
\end{tabular}
\caption{Kompaktlık varyantları}
\end{table}

\textbf{Sıralama:} Kompakt $\Rightarrow$ Sayılabilir kompakt $\Rightarrow$ Ardışımlı kompakt $\Rightarrow$ Lindelöf.

%-------------------------------------------------------------------
\section{Teoremler}

\begin{theorem}[Heine-Borel]
$\mathbb{R}^n$'de kompakt $\Leftrightarrow$ kapalı ve sınırlı.
\end{theorem}

\begin{theorem}
Kompakt $+$ Hausdorff $\Rightarrow$ Normal (T4).
\end{theorem}

\begin{theorem}
Kompakt uzaydan Hausdorff uzaya sürekli bijeksiyon $\Rightarrow$ homeomorfizma.
\end{theorem}

\begin{theorem}[Tychonoff]
Keyfi kompakt uzayların kartezyen çarpımı kompakttır.
\end{theorem}

\begin{theorem}[Alexandroff Tek-Nokta Kompaktifikasyonu]
Her lokal kompakt Hausdorff uzay $X$ için $X^* = X \cup \{\infty\}$ kompakt Hausdorff'tur.
\end{theorem}

%-------------------------------------------------------------------
\section{Algoritmalar}

\subsection{Sonlu Uzayda Kompaktlık}

Sonlu uzay her zaman kompakttır: her örtü zaten sonlu.
Algoritma: $O(1)$.

\subsection{analyze\_compactness\_variants — Tag Tabanlı Çıkarım}

\begin{algorithm}
\caption{AnalyzeVariants$(X, \tau)$}
\begin{algorithmic}[1]
\If{$X$ sonlu}
    \State tüm varyantlar $\leftarrow$ True; \Return VariantProfile(...)
\EndIf
\If{\texttt{'compact'} $\in$ tags}
    \State countably\_compact, sequentially\_compact, lindelof $\leftarrow$ True
\EndIf
\If{\texttt{'metric'} $\wedge$ \texttt{'compact'} $\in$ tags}
    \State sequentially\_compact $\leftarrow$ True
\EndIf
\end{algorithmic}
\end{algorithm}

%-------------------------------------------------------------------
\section{pytop API}

\begin{lstlisting}[language=Python]
from pytop import (
    is_compact,
    is_lindelof,
    is_locally_compact,
)
from pytop.compactness_variants import (
    is_countably_compact,
    is_sequentially_compact,
    analyze_compactness_variants,
)
\end{lstlisting}

\texttt{analyze\_compactness\_variants(space)} $\to$ \texttt{Result}; \texttt{.value} bir \texttt{dict} içerir.

%-------------------------------------------------------------------
\section{Örnekler}

\subsection*{Örnek 5.1 — Sonlu Uzaylar: Otomatik Kompakt}

\begin{lstlisting}[language=Python]
from pytop import sierpinski_space, discrete_topology, finite_chain_space, is_compact

s = sierpinski_space()
print("Sierpinski compact?", is_compact(s).status)
d = discrete_topology(1, 2, 3)
print("Discrete(3) compact?", is_compact(d).status)
c = finite_chain_space(3)
print("Chain(3) compact?", is_compact(c).status)
\end{lstlisting}

\begin{verbatim}
Sierpinski compact?      true
Chain(3) compact?        true
Discrete(3) compact?     true
\end{verbatim}

\subsection*{Örnek 5.2 — Gerçek Doğru $\mathbb{R}$: Kompakt Değil}

\begin{lstlisting}[language=Python]
rl = real_line_metric()
print("compact?", is_compact(rl).status)
print("lindelof?", is_lindelof(rl).status)
print("locally_compact?", is_locally_compact(rl).status)
\end{lstlisting}

\begin{verbatim}
compact?          false
lindelof?         true
locally_compact?  conditional
\end{verbatim}

\subsection*{Örnek 5.3 — analyze\_compactness\_variants: Sierpiński}

\begin{lstlisting}[language=Python]
r = analyze_compactness_variants(s)
for key, val in r.value.items():
    if hasattr(val, 'status'):
        print(f"  {key:<25}: {val.status}")
\end{lstlisting}

\begin{verbatim}
  representation           : finite
  countably_compact        : true
  sequentially_compact     : true
  pseudocompact            : true
  lindelof                 : true
\end{verbatim}

%-------------------------------------------------------------------
\section{Alıştırmalar}

\subsection*{Kodlama}

\begin{enumerate}[label=K\arabic*.]
\item \texttt{finite\_chain\_space(5)} ve \texttt{discrete\_topology(1,2,3,4,5)} için
      \texttt{is\_compact} karşılaştırın.
\item \texttt{naturals\_cofinite()} için \texttt{is\_compact}, \texttt{is\_lindelof},
      \texttt{is\_countably\_compact} hesaplayın.
\item \texttt{closed\_unit\_interval\_metric()} ile \texttt{real\_line\_metric()} için
      \texttt{analyze\_compactness\_variants} çalıştırıp karşılaştırın.
\end{enumerate}

\subsection*{Teori}

\begin{enumerate}[label=T\arabic*.]
\item Kompakt + sürekli $\Rightarrow$ görüntü kompakt olduğunu ispatlayın.
\item Heine-Borel teoreminin neden geçerli olduğunu: $[0,1]$ neden kompakt,
      $(0,1)$ neden kompakt değil?
\end{enumerate}
